\chapter{Proof of a Result about Homotopy Groups of CW-Pairs}

In this section, we give a detailed proof of the following proposition, which was used in the proof of Theorem \ref{psc:thm:existence_result_spin} and Theorem \ref{psc:thm:existence_result_totally_non_spin}.

\begin{proposition}
    \label{appendix2:prop:finite_generation_lemma}
    Let $(X,A)$ be a CW pair with $X$ connected. Suppose that $(X,A)$ is $(n-1)$-connected for $n \geq 2$, that the inclusion $A \to X$ induces an isomorphism on $\pi_1$, and that $(X,A)$ contains only finitely many $n$-cells. Then $\pi_n(X,A)$ is finitely generated as a $\Z[\pi_1(A)]$-module.
\end{proposition}

We will see in the proof that, even for $n=2$, the group $\pi_2(X,A)$ is always abelian, allowing us to view it as a $\Z[\pi_1(A)]$-module. Proposition \ref{appendix2:prop:finite_generation_lemma} rests on the following purely algebraic result, which is an adaptation of Lemma 1.12 in \cite{Pedersen2017}.

\begin{lemma}
    \label{appendix2:lem:homological_algebra_lemma}
    Let $P_\ast$ be a positive chain complex of projective $R$-modules, which is homologically $(n-1)$-connected, that is, $P_i = 0$ for $i < 0$ and $H_k(P_\ast) = 0$ for $k \leq n-1$. Suppose that $P_n$ is finitely generated. Then $H_n(P_\ast)$ is also finitely generated.
\end{lemma}

\begin{proof}
    Denote by $\partial_k \colon P_k \to P_{k-1}$ the $k$-th boundary map. We show inductively that $P_k \cong \mathrm{ker}(\partial_k) \oplus \mathrm{ker}(\partial_{k-1})$ for $0 \leq k \leq n$. The base case $k = 0$ is trivially true, so suppose the claim holds for some $k < n$. Then $\mathrm{ker}(\partial_k)$ is projective as a direct summand of a projective module, and $H_{k}(P_\ast) = 0$ is equivalent to the fact that $\partial_{k+1} \colon P_{k+1} \to \mathrm{ker}(\partial_k)$ is surjective. Hence $P_{k+1} \cong \mathrm{ker}(\partial_{k+1}) \oplus \mathrm{ker}(\partial_k)$. This completes the induction. So, in particular,
    \[
    P_n \cong \mathrm{ker}(\partial_n) \oplus \mathrm{ker}(\partial_{n-1}).
    \]
    Therefore, $\mathrm{ker}(\partial_n)$ is finitely generated as a direct summand of a finitely generated module. Consequently, the quotient $H_n(P_\ast)$ is also finitely generated. 
\end{proof}

Moreover, we need the following fact concerning the homology of relative CW-complexes.

\begin{lemma}
    \label{appendix2:lem:finite_generated_homology}
Let $Y$ and $B$ be topological spaces, where $Y$ is obtained from $B$ by attaching finitely many $n$-cells. Let $p \colon \hat{Y} \to Y$ be some cover, such that the group $\mathrm{Deck}(p)$ of deck transformations of $p$ acts transitively on the fibers. Set $\hat{B} := p^{-1}(B)$. Then $H_n(\hat{Y},\hat{B})$ is finitely generated as a $\Z[\mathrm{Deck}(p)]$-module, where the action of an element $\phi \in \mathrm{Deck}(p)$ is given by $\phi_\ast \colon H_n(\hat{Y},\hat{B}) \to H_n(\hat{Y},\hat{B})$.
\end{lemma}

\begin{proof}
Let us recall how the relative CW-structure of $(Y,B)$ lifts to a relative CW-structure of $(\hat{Y},\hat{B})$. Let $I$ be the index set for the $n$-cells of $(Y,B)$ with characteristic maps $Q_i \colon D^n \to Y$, $i \in I$. Choose some fixed basepoint $e_0 \in D^n$. Then for each $y \in \tilde{Y}$ with $p(y) = Q_i(e_0)$, there exists a unique lift $Q_i^y \colon D^n \to \hat{Y}$ of $Q_i$ along $p$, such that $Q_i^y(e_0) = y$. These $Q_i^y$, indexed over the set
$\Lambda := \{(y,i) \in \hat{Y} \times I \mid p(y) = Q_i(e_0)\}$,
are then the characteristic maps for the $n$-cells of $(\hat{Y},\hat{B})$. The corresponding attaching maps are the restrictions $q_i^y := Q_i^y|_{S^{n-1}} \colon S^{n-1} \to \hat{B}$.
See \cite[Theorem 1.1.6]{fritsch1990cellular} for further details. From excision in the pushout

\[\begin{tikzcd}
{\coprod_{\Lambda}S^{n-1}} & {\hat{B}} \\
{\coprod_{\Lambda}D^n} & {\hat{Y}}
\arrow["{(q_i^y)}", from=1-1, to=1-2]
\arrow[from=1-1, to=2-1]
\arrow[from=1-2, to=2-2]
\arrow["{(Q_i^y)}", from=2-1, to=2-2]
\end{tikzcd}\]

\noindent we obtain an isomorphism
\[
\bm{Q} \colon \bigoplus_{\Lambda} H_n(D^n,S^{n-1}) \to H_n(\hat{Y},\hat{B}).
\]

An element $\phi \in \mathrm{Deck}(p)$ acts on $\bigoplus_{\Lambda} H_n(D^n,S^{n-1})$ by mapping the summand with index $(y,i)$ to the summand with index $(\phi(y),i)$. We claim that $\bm{Q}$ is equivariant with respect to this action and the previously defined $\mathrm{Deck}(p)$-action on $H_n(\hat{Y},\hat{B})$. This follows from commutativity of the following diagram, where again the action of an element $\phi \in \mathrm{Deck}(p)$ on $\coprod_{\Lambda} D^n$ and $\coprod_{\Lambda} S^{n-1}$ is given by permuting indices as above:

\[
\begin{tikzcd}[row sep={40,between origins},column sep={60,between origins},every label/.append style={font=\small}]
&\coprod_{\Lambda}S^{n-1}\ar["(q_i^y)"{pos=.55}]{rr}\ar["b"{pos=.5}]{dd}&&\hat{B}\ar{dd}\\
\coprod_{\Lambda}S^{n-1}\ar["(q_i^y)"{pos=.55},crossing over]{rr}\ar{dd}\ar["\phi"{pos=.5}]{ur}&&\hat{B}\ar["\phi"{pos=.5}]{ur}\\
&\coprod_{\Lambda}D^n\ar["(Q_i^y)"{pos=.275}]{rr}&&\hat{Y}\\
\coprod_{\Lambda}D^n\ar["(Q_i^y)"{pos=.4}]{rr}\ar["\phi"{pos=.5}]{ur}&&\hat{Y}\ar[from=uu,crossing over]\ar["\phi"{pos=.5}]{ur}
\end{tikzcd}
\]  

Therefore, $\bm{Q}$ is an isomorphism of $\Z[\mathrm{Deck}(p)]$-modules. But since $\mathrm{Deck}(p)$ acts transitively on the fibers of $p$ and we assumed the index set $I$ to be finite, $\bigoplus_{\Lambda} H_n(D^n,S^{n-1})$ is a finitely generated $\Z[\mathrm{Deck}(p)]$-module (for each $i \in I$, choose one $(y,i) \in \Lambda$ and a generator of the summand $H_n(D^n,S^{n-1})$ corresponding to this index $(y,i)$).
\end{proof}

\begin{proof}[Proof of Proposition \ref{appendix2:prop:finite_generation_lemma}]
    Let $p \colon \tX \to X$ be the universal cover. Since the inclusion $A \to X$ induces an isomorphism on $\pi_1$, the universal cover of $A$ is given by $\tA = p^{-1}(A)$. Choosing basepoints $x_0 \in A$ and $\tx_0 \in \tA$ with $p(\tx_0) = x_0$, we obtain canonical isomorphisms 
    \[
    \pi_1(A,x_0) \cong \pi_1(X,x_0) \cong \mathrm{Deck}(p),
    \]
    where $\mathrm{Deck}(p)$ is the group of deck transformations of $p$. By virtue of these isomorphisms, we denote any of the three groups simply by $\pi$. Consider the following actions of $\pi$. On $\pi_n(X,A,x_0)$ we have the natural $\pi_1(A,x_0)$-action. On $H_n(\tX,\tA)$ consider the action
    \[
    \mathrm{Deck}(p) \ni \phi \mapsto \phi_\ast \colon H_n(\tX,\tA) \to H_n(\tX,\tA)
    \]
    and on $\pi_n(\tX,\tA,\tx_0)$ the action 
    \[
    \mathrm{Deck}(p) \ni \phi \mapsto \phi_\ast \colon \pi_n(\tX,\tA,\tx_0) \to \pi_n(\tX,\tA,\phi(\tx_0)) \cong \pi_n(\tX,\tA,\tx_0),
    \]
    where the second isomorphism is the canonical change of basepoint isomorphism ($\tA$ is simply connected, so this isomorphism does not depend on the chosen path from $\tx_0$ to $\phi(\tx_0)$). Consider the isomorphisms
    \[
    p_\ast \colon \pi_n(\tX,\tA,\tx_0) \to \pi_n(X,A,x_0)
    \]
    and 
    \[
    h \colon \pi_n(\tX,\tA,\tx_0) \to H_n(\tX,\tA),
    \]
    where $h$ is the Hurewicz map. Note that $h$ is an isomorphism by the Hurewicz Theorem since $\pi_k(\tX,\tA,\tx_0) \cong \pi_k(X,A,x_0) = 0$ for $2 \leq k \leq n-1$, and $\tX, \tA$ are simply connected. These isomorphisms are now equivariant with respect to the $\pi$-actions defined above. For $p_\ast$, this is a direct consequence of the definitions, and for $h$, this follows from naturality of the Hurewicz map and the fact that $h([f])$ only depends on the free homotopy class of $f \colon (D^n,S^{n-1}) \to (\tX,\tA)$, which is invariant under the change of basepoint isomorphism. So, in other words, $p_\ast$ and $h$ are isomorphisms of $\Z[\pi]$-modules (note that even for $n=2$ all groups are abelian). Therefore, it suffices to show that $H_n(\tX,\tA)$ is a finitely generated $\Z[\pi]$-module. Let $C_\ast^\mathrm{cell}(\tX,\tA)$ be the cellular chain complex of $(\tX,\tA)$ with respect to the relative CW-structure induced by $(X,A)$. If $X_k$ denotes the $k$-skeleton of $(X,A)$, the $k$-skeleton of $(\tX,\tA)$ is given by $\tX_k = p^{-1}(X_k)$. Compare the proof of Lemma \ref{appendix2:lem:finite_generated_homology} or \cite[Proposition 2.3.9]{fritsch1990cellular}. By assumption, $X_n$ is obtained from $X_{n-1}$ by attaching finitely many $n$-cells. Consider the cover $p|_{\tX_n} \colon \tX_n \to X_n$ and note that $\pi = \mathrm{Deck}(p) \cong \mathrm{Deck}(p|_{\tX_n})$, the isomorphism being given by restricting a deck transformation of $p$ to the $n$-skeleton $\tX_n$. Clearly $\mathrm{Deck}(p|_{\tX_n})$ still acts transitively on the fibers. Therefore, we can apply Lemma \ref{appendix2:lem:finite_generated_homology} to conclude that 
    \[
    C_n^\mathrm{cell}(\tX,\tA) = H_n(\tX_{n},\tX_{n-1})
    \]
    is a finitely generated $\Z[\pi]$-module. Moreover, it follows from the exact sequence
    \[\begin{tikzcd}[column sep = 1em]
	{H_1(\tilde{X})} & {H_1(\tilde{X}, \tilde{A})} & {H_0(\tilde{A})} & {H_0(\tilde{X})} & {H_0(\tilde{X}, \tilde{A})}
	\arrow[from=1-1, to=1-2]
	\arrow[from=1-2, to=1-3]
	\arrow[from=1-3, to=1-4]
	\arrow[from=1-4, to=1-5]
\end{tikzcd}\]
that $H_0(\tilde{X}, \tilde{A}) = H_1(\tilde{X}, \tilde{A}) = 0$, since both $\tilde{X}$ and $\tilde{A}$ are connected and $\tilde{X}$ is simply connected. Furthermore, 
\[
H_k(\tilde{X}, \tilde{A}) \cong \pi_k(\tilde{X}, \tilde{A}, \tx_0) \cong \pi_k(X,A,x_0) = 0
\]
for $2 \leq k \leq n-1$. Note that, by naturality, the boundary maps of the chain complex $C_\ast^\mathrm{cell}(\tX,\tA)$ are also $\Z[\pi]$-linear. Hence $C_\ast^\mathrm{cell}(\tX,\tA)$ satisfies the conditions of Lemma \ref{appendix2:lem:homological_algebra_lemma}, and we conclude that $H_n(\tilde{X}, \tilde{A})$ is indeed a finitely generated $\Z[\pi]$-module.
\end{proof}